\subsection{Diseño experimental}

La campaña experimental se definió a partir de cargas de trabajo explícitas para cada
circuito, con el fin de separar escenarios con alta regularidad de escenarios adversos para
la compresión del vector de estado. En todos los casos se evaluaron 22, 23 y 24 qubits. Cada
comparación entre estrategias se realizó sobre la misma instancia del circuito, con los
mismos parámetros de entrada y bajo la misma política de recolección de métricas.

En Grover se analizaron dos casos. El primero correspondió a un oráculo con objetivo único y
se ejecutó con tres posiciones representativas del estado marcado: $N/8$, $N/2$ y $7N/8$,
donde $N=2^n$ es el tamaño del espacio de búsqueda. El segundo correspondió a un oráculo
multiobjetivo con tres estados marcados ubicados en esas mismas regiones. En QFT se
consideraron dos familias de entrada: una superposición periódica y una superposición de alta
entropía con amplitudes de magnitud uniforme y fases aleatorias,
$a_y = e^{i\phi_y}/\sqrt{N}$ con $\phi_y \sim U[0, 2\pi)$ y semilla fija.

Las ejecuciones se realizaron en la estación de trabajo utilizada para el desarrollo del simulador: Rocky Linux 10.1, procesador AMD Ryzen 7 5700X de 8 núcleos y 16 hilos, y 32~GiB de memoria RAM. La campaña se instrumentó con los scripts de perfilado del repositorio del simulador, que ejecutan los binarios experimentales y recolectan tiempo transcurrido y memoria residente máxima mediante \texttt{psrecord}. En Grover, cada fila resume tres instancias por tamaño y estrategia; en QFT, cada fila corresponde a una instancia por familia de entrada, número de qubits y estrategia.

\begin{table}[H]
    \centering
    \small
    \caption{Cargas de trabajo utilizadas en la evaluación}
    \label{tab:exp-workloads}
    \begin{tabular}{>{\raggedright\arraybackslash}p{2.1cm}>{\raggedright\arraybackslash}p{2.8cm}>{\centering\arraybackslash}p{1.7cm}>{\raggedright\arraybackslash}p{6.8cm}}
        \hline
        \textbf{Circuito} & \textbf{Caso} & \textbf{Qubits} & \textbf{Definición de la entrada} \\
        \hline
        Grover & Objetivo único & 22, 23 y 24 & Tres instancias por tamaño con estado marcado en $N/8$, $N/2$ y $7N/8$. \\
        Grover & Multiobjetivo & 22, 23 y 24 & Tres estados marcados por tamaño, ubicados en $N/8$, $N/2$ y $7N/8$. \\
        QFT & Superposición periódica & 22, 23 y 24 & Estado de entrada con patrón regular y repetitivo. \\
        QFT & Superposición de alta entropía & 22, 23 y 24 & Estado de entrada con magnitud uniforme y fases aleatorias. \\
        \hline
    \end{tabular}
\end{table}

Las estrategias comprimidas se evaluaron con configuraciones efectivas fijas, resumidas en la
Tabla~\ref{tab:exp-configs}. Estas configuraciones corresponden a los parámetros finalmente
resueltos por el simulador para las cargas de trabajo Grover y QFT, bien sea por valores por
defecto o por la sintonía automática guiada por el tipo de acceso al estado. La estrategia de referencia corresponde al almacenamiento no
comprimido del simulador y sirve como base para calcular las razones de memoria y tiempo
reportadas en las tablas de resultados.

\begin{table}[H]
    \centering
    \small
    \caption{Configuraciones de compresión utilizadas en toda la campaña experimental}
    \label{tab:exp-configs}
    \begin{tabular}{>{\raggedright\arraybackslash}p{2.1cm}>{\raggedright\arraybackslash}p{9.2cm}}
        \hline
        \textbf{Estrategia} & \textbf{Configuración} \\
        \hline
        Blosc & $\texttt{chunkStates}=32768$, $\texttt{gateCacheSlots}=8$, $\texttt{clevel}=1$, $\texttt{nthreads}=4$, $\texttt{compcode}=1$, $\texttt{useShuffle}=\texttt{true}$. \\
        ZFP & Modo $\texttt{FixedPrecision}$, $\texttt{precision}=40$, $\texttt{chunkStates}=32768$, $\texttt{gateCacheSlots}=8$, $\texttt{nthreads}=4$. \\
        \hline
    \end{tabular}
\end{table}

Las métricas principales fueron el tiempo total de ejecución y la memoria residente máxima.
Ambas se obtuvieron a partir de las trazas \texttt{psrecord}: el tiempo corresponde al último
valor registrado en la columna \textit{Elapsed time} y la memoria máxima al máximo de la
columna \textit{Real (MB)}. A partir de estas mediciones se definieron dos indicadores de
lectura directa:
\[
\text{Ahorro de memoria (\%)} = 100\left(1-\frac{M_{\text{comp}}}{M_{\text{ref}}}\right),
\qquad
\text{Tiempo relativo (x)} = \frac{T_{\text{comp}}}{T_{\text{ref}}}.
\]
En estas expresiones, $M_{\text{ref}}$ y $T_{\text{ref}}$ corresponden a la ejecución con
almacenamiento no comprimido del mismo circuito y del mismo número de qubits. Un ahorro de
memoria negativo indica que la estrategia evaluada consumió más memoria que la referencia.
Para Blosc, la exactitud se verificó exportando el vector completo de amplitudes de la
ejecución comprimida y comparándolo amplitud por amplitud contra la referencia no comprimida de la
misma instancia experimental; en todos los casos reportados esa comparación produjo igualdad
exacta y, por consiguiente, error máximo absoluto nulo. Para ZFP se deja el
campo de error explícito en las tablas con el fin de consignar el valor medido correspondiente
mediante la misma comparación completa del estado.
En la comparación con la variante optimizada de Grover no se reporta error, porque allí el
interés ya no es contrastar exactitud entre estrategias de compresión, sino cuantificar cómo
cambia la relación entre memoria y tiempo cuando la simulación explota directamente la
estructura algebraica del algoritmo.
